Šajā nodaļā izklāstīta promocijas darba pētījumos un to rezultātā tapušajās publikācijās izstrādātā realizācijas metodoloģija. Tā neapraksta tikai vienu sistēmu, bet apvieno vairāku publikāciju ieguldījumu kopīgā neiromorfās skaitļošanas ietvarā.

Darbs aptver realizācijas no impulsu neironu modeļiem līdz pilnīgai neiromorfās uztveres sistēmai. Vispirms pētītas bioloģiski iedvesmotu neironu modeļu aparatūras realizācijas. Tās izmantotas kā uz notikumiem balstītas apstrādes pamatelementi. Uz to pamata izstrādāts pārkonfigurējams \Gls{fpga} paātrinātājs konvolucionālo \Gls{snn} paātrināšanai aparatūrā. Visbeidzot izstrādāti dziļie konvolucionālie impulsu neironu tīkli praktiskiem uztveres uzdevumiem, izmantojot retinātus 3D sensoru datus.

Pārējā nodaļa sadalīta divās saistītās daļās. Pirmajā uzmanība pievērsta PCIe \Gls{fpga} paātrinātāja aparatūras arhitektūrai pārkonfigurējamu konvolucionālo \Gls{snn} paātrināšanai gan ar \Gls{lif}, gan Ižikeviča neironu modeļiem. Otrajā demonstrēta dziļo \Gls{snn} izmantošana 3D objektu detektēšanā ar LiDAR un lidojuma laika sensoriem, parādot modeļu un aparatūras risinājuma izmantošanu praktiskās autonomās uztveres sistēmās.

Realizācijas balstās uz promocijas darba izstrādes laikā publicētajiem pētījumiem. Ižikeviča neirona aparatūras realizācija pirmoreiz publicēta darbā \cite{lielamurs2023a}. Retināta impulsu objektu detektēšana ar LiDAR datiem balstās uz darbu \cite{lielamurs2024}, savukārt tās paplašinājums automobiļu \Gls{tof} sensoru sistēmai publicēts darbā \cite{lielamurs2025}. Visbeidzot pārkonfigurējamā \Gls{fpga} paātrinātāja arhitektūra ir darba \cite{lielamurs2026} galvenais ieguldījums.


